\section{Critical discussion and limitations}
The strongest result is the agreement between an exact mathematical prediction and every controlled E1 subset observation. It depends on the same integer $M$, valid textbook RSA inputs, $e=3$, and pairwise-coprime moduli. The leading-zero comparison exposes an important boundary of this statement: modular arithmetic concerns integer representatives, whereas an application transports byte strings with an encoding convention. A demonstration that visually prints the original 96 bytes must disclose that its evaluator supplied the missing length. This distinction also prevents an exact cube root from being mistaken for authenticated knowledge of a sender's original data.

E2 improves causal interpretation by changing the encryption scheme within a pair while holding the keys and message fixed. Its 0/60 OAEP recovery count is still evidence about one selected CRT-and-root procedure and three message samples, not a statistical or formal proof of OAEP security. The 60 comparisons are correlated within 20 independently generated key-set blocks. Likewise, ciphertext diversity under repeated OAEP encryption illustrates its randomized encoding but does not replace a security argument in an explicit adversarial model. The project uses established OAEP through a maintained cryptographic library rather than implementing padding itself.

E3 distinguishes cryptographic computation from application packaging. The dominant 1\,MiB serialization interval suggests that a more compact representation would be a sensible \emph{future engineering comparison}, but the present data do not measure such an alternative. JSON/Base64 improves inspectability and portability at a measurable space and time cost. E3's timing includes Python and operating-system effects on one Windows computer; its 20 repetitions per configuration describe this run, not every processor or cloud environment. Recipient-key generation is separate, and one-recipient decryption timing cannot predict the cost of sequentially decrypting for all recipients.

Secure sharing assumes each recipient public key has been authenticated before encryption. The package's fingerprint is a stable identifier, not proof of personal identity. AES-GCM verifies consistency of ciphertext and AAD under the unwrapped key~\cite{dworkin2007gcm}; it does not establish who sent the package. The current system has no sender signature, replay protection, certificate infrastructure, revocation or persistent private-key management. Its visible filename and recipient metadata are not confidential. E4's selected rejections neither cover all malformed inputs nor distinguish accidents from adversarial modifications. Extending the system would require an explicit threat model and new evidence, not simply more interface controls.

\section{Conclusion}
BroadcastGuard links a public-only RSA broadcast attack to its exact numerical recovery condition, a same-key randomized-padding comparison and a functioning authenticated multi-recipient file package. The 560 E1 observations confirm the known CRT bound for the declared inputs while showing why byte count and integer value differ. The 60 E2 pairs demonstrate non-recovery by the tested method for OAEP ciphertexts alongside successful recipient decryption. E3 reveals a concrete JSON/Base64 packaging cost, and E4 verifies selected fail-closed outcomes. The work's contribution is the controlled integration, validation and interpretation of established methods. It does not introduce a new cryptographic primitive or establish universal security.
